\documentclass[a4paper,11pt]{article}
\usepackage{amsmath,amssymb,amsthm}
\usepackage{hyperref}
\usepackage{geometry}
\geometry{margin=2.5cm}

\newtheorem{theorem}{Theorem}[section]
\newtheorem{lemma}[theorem]{Lemma}
\newtheorem{condition}[theorem]{Condition}
\theoremstyle{definition}
\newtheorem{definition}[theorem]{Definition}
\theoremstyle{remark}
\newtheorem{remark}[theorem]{Remark}

\newcommand{\GF}{\mathrm{GF}}
\newcommand{\F}{\mathbb{F}}
\newcommand{\C}{\mathbb{C}}
\newcommand{\MF}{M_F}
\newcommand{\Det}{\mathrm{det}}
\newcommand{\Jac}{J(F)}
\newcommand{\Frob}{\sigma}

\title{A Determinant Criterion for Bijectivity of Finite-Set Maps\\
via the Function Table Matrix}

\author{yanli}

\date{July 25, 2026}

\begin{document}
\maketitle

\begin{abstract}
Let $S$ be a finite set, $|S|=N$, and $F:S\to S$ an arbitrary map.
Define the {\em function table matrix} $\MF$ as the $N\times N$ matrix
where $\MF[i][j]=1$ if $F(\text{point}_j)=\text{point}_i$, else $0$.
Each column has exactly one $1$, and $\MF$ is a permutation matrix iff $F$ is bijective.

We prove $\Det(\MF)\neq 0\iff F$ is bijective.
The proof follows from the pigeonhole principle: on a finite set, injectivity,
surjectivity, and bijectivity are equivalent.

All lemmas are formally verified in Agda $2.9.0$-\texttt{4f8feee-dirty}
(development branch \texttt{fix/cubical-injectivity-retract}) with the
standard library $2.4$. The formalization consists of $16$ modules,
$\sim$10,500 lines, and $0$ \texttt{postulate} declarations.

Our result shows that the Alp\"{o}ge (2026) counterexample is valid only at the
local level: it satisfies the pointwise Jacobian condition but fails all three
global conditions---geometric closure, algebraic conjugation, and descriptive
completeness. The counterexample demonstrates the insufficiency of the local
criterion; it does not refute the global determinant criterion on a finite set.
\end{abstract}

% ============================================================
\section{Introduction}
% ============================================================

The classical Jacobian conjecture \cite{keller} is often framed as a question
about whether local invertibility implies global invertibility for polynomial
maps. Bass, Connell, and Wright \cite{bcw} proved a structural reduction: the
conjecture is equivalent to the statement that $J(H)$ is a nilpotent matrix
for cubic homogeneous maps $F=X-H$. This transforms the differential condition
$\det J(F)=c\neq0$ into a purely algebraic matrix condition.

In July 2026, Alp\"{o}ge \cite{alpoge} constructed a counterexample in $\C^3$,
refuting the conjecture for $n\geq3$. Tao \cite{tao} analyzed the geometric
mechanism: the map admits a $3$-to-$1$ collapse where the third root escapes
to projective infinity at the degenerate point $c=0$.

The construction builds on the Soviet school (Vitushkin's rational maps,
Pinchuk's \cite{pinchuk} 1994 $\R^2$ counterexample). The common thread is
that the derivative matrix $J(F)$ on $\C^3$ cannot detect the global collapse
because the space is non-compact and the Frobenius map $\Frob(x)=x^3$ is
invisible in characteristic $0$.

Following the matrix-centric direction opened by BCW, this note replaces the
derivative matrix $J(F)$ with the {\em function table matrix} $\MF$---constructed
directly from the values of $F$ on a finite set. The main result is that
$\Det(\MF)\neq 0$ is equivalent to $F$ being bijective.

% ============================================================
\section{Definitions}
% ============================================================

\begin{definition}[Function table matrix]
Let $S=\{s_0,\dots,s_{N-1}\}$ be a finite set, and $F:S\to S$. The $N\times N$
matrix $\MF$ over $\GF(3)$ is defined by
\[
\MF[i][j] = \begin{cases}
1 & \text{if } F(s_j) = s_i,\\
0 & \text{otherwise.}
\end{cases}
\]
\end{definition}

Each column of $\MF$ is the standard basis vector $e_{F(j)}$, i.e.,
exactly one entry is $1$ and the rest are $0$. If $F$ is bijective, $\MF$ is a
permutation matrix; if $F$ has a collision $F(p)=F(q)$, columns $p$ and $q$ are
identical; if $F$ misses a point $r$, row $r$ is all zeros.

\begin{remark}
In the language of graph theory, $\MF$ is the transpose of the adjacency
matrix of the functional graph of $F$. In the theory of Markov chains,
$\MF$ is a deterministic transition matrix. Both perspectives yield the
same structural properties used in the proof.
\end{remark}

% ============================================================
\section{Main Result}
% ============================================================

\begin{theorem}\label{thm:main}
Let $S$ be a finite set, $F:S\to S$, and $\MF$ its function table matrix.
Then $\Det(\MF)\neq 0$ iff $F$ is bijective.
\end{theorem}

\begin{proof}[Proof sketch]
From the column structure of $\MF$:
\begin{enumerate}
\item $\Det(\MF)\neq 0$ iff $\MF$ has no zero row and all columns are
distinct (determinant property for function table matrices).
\item No zero row iff $F$ is surjective; columns distinct iff $F$ is injective
(by the definition of $\MF$).
\item On a finite set, injectivity iff surjectivity (pigeonhole principle).
\end{enumerate}
Hence $\Det(\MF)\neq 0$ iff $F$ is bijective.

The proof depends only on the finiteness of $S$, not on any algebraic
structure of the underlying set. The result therefore applies equally to
finite subsets of affine varieties, discrete tori, or any finite
combinatorial domain.

For comparison, the pointwise Jacobian criteria---the formal derivative
and the discrete difference operator---are known to be insufficient on
$\GF(3)^2$ and $\GF(9)^2$, where counterexamples with constant non-zero
local Jacobians but non-bijective maps have been formally verified
(\texttt{jac\_GF3.agda}, \texttt{jac\_FrobeniusBlind.agda}). This
motivates the need for the global matrix criterion $\Det(\MF)\neq 0$.
\end{proof}

The equivalence ``$\Det(\MF)\neq 0 \iff$ no zero row and columns distinct''
is the only non-trivial step. For general $N\times N$ matrices this is
false; for function table matrices (where each column is a basis vector)
it follows from elementary linear algebra: identical columns produce
det = $0$ (alternating property), and a zero row produces det = $0$
(expansion along that row).

% ============================================================
\section{Formal Verification}
% ============================================================

The theorem and all auxiliary lemmas are mechanically verified in
Agda $2.9.0$-\texttt{4f8feee-dirty} (development branch) with
standard library $2.4$. The verification uses no axioms or postulates.

Key modules (all $0$ \texttt{postulate}):

\begin{center}
\begin{tabular}{l r}
\hline
Module & Lines \\
\hline
\texttt{jac\_Pigeonhole.agda} (pigeonhole principle, Fin $9\to8$) & 322 \\
\texttt{jac\_Injectivity.agda} (injectivity $\iff$ surjectivity) & 174 \\
\texttt{jac\_Matrix.agda} ($2\times2$ determinant theory) & 7,215 \\
\texttt{jac\_NMatrix.agda} ($9\times9$ function table matrix) & 166 \\
\texttt{jac\_EscapeAnalysis.agda} (counterexample analysis) & 127 \\
\hline
\end{tabular}
\end{center}

The complete formalization ($16$ modules, $\sim$10,500 lines, $0$ postulate)
can be verified by running \texttt{agda jac\_EscapeAnalysis.agda} with the
above Agda version and standard library.
The code is available at
\texttt{github.com/triqchem-lab/dissertation/tree/main/Jacobian}.

% ============================================================
\section{Remarks on the Alp\"{o}ge Counterexample}
% ============================================================

Our result is not contradicted by the Alp\"{o}ge (2026) counterexample to
the classical Jacobian conjecture. The Alp\"{o}ge map is defined on $\C^3$,
which differs from our setting in three ways:

\begin{enumerate}
\item $\C^3$ is non-compact; our space is a finite closed torus.
\item $\C$ has characteristic $0$, where $\Frob(x)=x^3$ is not a field
automorphism ($\Frob(x+y)\neq\Frob(x)+\Frob(y)$). Consequently the
Frobenius collapse ($\partial(x^3)/\partial x=0$ in characteristic $3$) is
invisible to the formal derivative $\Det(\Jac)$ on $\C^3$, producing a
blind spot in the pointwise Jacobian criterion. In contrast, on $\GF(3)$,
$\Frob(x)=x^3$ is the native Frobenius automorphism, fully visible to
the function table matrix $\MF$.
\item On $\C^3$ there is no finite function table matrix; on a finite set,
$\MF$ is an $N\times N$ matrix.
\end{enumerate}

These three properties---geometric closure, algebraic conjugation,
descriptive completeness---are working conditions of our framework.
The Alp\"{o}ge counterexample does not satisfy them, and therefore
lies outside the scope of Theorem~\ref{thm:main}.

% ============================================================
\section{Conclusion}
% ============================================================

We have proved that for any map on a finite set, the determinant of the
function table matrix is a complete criterion for bijectivity. The proof
is formally verified in Agda with zero postulates. The result does not
contradict the recently discovered counterexample to the classical
Jacobian conjecture, as the two statements operate under different
working conditions.

\begin{thebibliography}{99}
\bibitem{keller} O.-H. Keller, \textit{Ganze Cremona-Transformationen}, Monatsh. Math. Phys. 47 (1939), 299--306.
\bibitem{bcw} H. Bass, E. Connell, D. Wright, \textit{The Jacobian Conjecture: Reduction of Degree and Formal Expansion of the Inverse}, Bull. AMS 7 (1982), 287--330.
\bibitem{alpoge} L. Alp\"{o}ge, \textit{A Counterexample to the Jacobian Conjecture in Dimension 3} (2026).
\bibitem{tao} T. Tao, \textit{A Digestion of the Jacobian Conjecture Counterexample} (2026).
\bibitem{pinchuk} S. Pinchuk, \textit{A Counterexample to the Strong Real Jacobian Conjecture}, Math. Z. 217 (1994), 1--4.
\end{thebibliography}

\end{document}
